%=====================================================================
\section{Admitting change}\label{sec:admission}

We now treat every finding as the outcome of a \emph{slot}: a pair of rules, a cross-layer cell $(u, t, r)$, a usage cell $(t, r, s)$ or $(t, r, \rho)$, or a drift cell of a tool. A state assigns each slot a finding or nothing. This section shows that an event can change only the slots of a small region, and that checking that region decides whether the invariant survives the event.

\subsection{Regions}\label{sec:region}
Let $\mathrm{rd}(r) = \{t : r \in \reads(d_1(t))\}$ be the tools that currently read $r$.

\begin{definition}[Region of an event]\label{def:region}
For a state $\G$ and an event $\epsilon$, the region $\Reg(\epsilon)$ is the following set of slots.
\begin{itemize}
\item $\mathsf{ins}(e)$ or $\mathsf{del}(e)$ with $e$ in layer $\ell$: the pairs $\{e, e'\}$ with $e'$ in layer $\ell$ and $\key(e) \cap \key(e') \neq \emptyset$, together with
\begin{itemize}
\item if $\ell = \mathsf{data}$: the cells $(u, t, r)$ with $u \in S$, $r \in \key(e)$, $t \in \mathrm{rd}(r)$;
\item if $\ell = \mathsf{tool}$: the cells $(u, t, r)$ with $u \in S$, $t \in \key(e)$, $r \in \reads(d_1(t))$;
\item if $\ell = \mathsf{usage}$: the usage and drift cells of the tools in $\mathrm{rd}(r)$, $r \in \key(e)$.
\end{itemize}
\item $\mathsf{redesc}(t, d)$: every cross-layer, usage and drift cell of $t$, for the resources of $\reads(d_1(t)) \cup \reads(d)$.
\item $\mathsf{approve}(t)$: the drift cells of $t$.
\end{itemize}
\end{definition}
The region is defined by key vertices and by the reads relation, both of which are indexed; it does not depend on the rest of the hypergraph.

\begin{theorem}[Locality]\label{thm:locality}
For every state $\G$ and event $\epsilon$, every slot outside $\Reg(\epsilon)$ has the same finding in $\G$ and in $\G \oplus \epsilon$. Consequently
\begin{equation}
\F(\G \oplus \epsilon) = \bigl(\F(\G) \setminus \Reg(\epsilon)\bigr) \cup \F_{\Reg(\epsilon)}(\G \oplus \epsilon), \label{eq:locality}
\end{equation}
where $\F_X$ denotes the findings in the slots of $X$.
\end{theorem}
\begin{proof}
\emph{Rule events.} Pairs that do not contain $e$ keep their match sets, effects and owners. A pair $\{e, e'\}$ with $\key(e) \cap \key(e') = \emptyset$ has $c(\kappa) \cap c'(\kappa) = \emptyset$ on the key attribute $\kappa$, so it neither overlaps (\Cref{lem:overlap}) nor stands in subsumption (\Cref{lem:subsume}, since $c(\kappa)$ is non-empty), and is not a finding before or after the event. By \Cref{def:rule}, $\dec_\ell(q)$ depends only on the rules matching $q$, and a rule matches a data or usage request on $r$ only if $r$ is among its key vertices, and a tool request on $t$ only if $t$ is. Hence a data event changes $\dec_{\mathsf{data}}(i, r)$ only for $r \in \key(e)$, a tool event changes $\dec_{\mathsf{tool}}(u, t)$ only for $t \in \key(e)$, and a usage event changes $\dec_{\mathsf{usage}}(r, p, s)$ only for $r \in \key(e)$. The cells of \Cref{def:cross} are conjunctions of $\dec_{\mathsf{tool}}(u, t)$ and $\dec_{\mathsf{data}}$ on the resource of the cell, so they change only in the cells listed. A usage cell of $t$ on $r$ reads $\dec_{\mathsf{usage}}$ on $r$ only. $\Drift_t$ contains only elements of $\Viol(d_1(t))$, whose resources are read by $t$; changes of $\Viol(d_0(t))$ on a resource that $d_1(t)$ does not read cannot remove any of them.
\emph{Tool events.} Every slot of \Cref{def:cross,def:usage} other than pairs is indexed by one tool and reads the description of that tool only, so $\mathsf{redesc}(t, d)$ changes only cells of $t$, and $\mathsf{approve}(t)$ only $\Drift_t$.
Equation~\eqref{eq:locality} restates the claim.
\end{proof}

\subsection{Admission}\label{sec:admit}
\begin{definition}[Admission check]\label{def:admission}
An event $\epsilon$ is \emph{admitted} in state $\G$ when $\F_{\Reg(\epsilon)}(\G \oplus \epsilon) \subseteq \Acc$. Otherwise it is \emph{held}, and the findings $\F_{\Reg(\epsilon)}(\G \oplus \epsilon) \setminus \Acc$ are sent to the owners they involve.
\end{definition}

\begin{theorem}[Admission is necessary and sufficient]\label{thm:preservation}
Let $\G$ be a governed state and $\epsilon$ an event. Then $\epsilon$ is admitted if and only if $\G \oplus \epsilon$ is governed. Consequently, if $\G_0$ is governed and only admitted events are applied, every reachable state is governed. If $\G_0$ is not governed, applying only admitted events never adds an unaccepted finding: $\F(\G \oplus \epsilon) \setminus \Acc \subseteq \F(\G) \setminus \Acc$.
\end{theorem}
\begin{proof}
By \eqref{eq:locality}, $\F(\G \oplus \epsilon) \setminus \Acc = \bigl((\F(\G) \setminus \Reg(\epsilon)) \setminus \Acc\bigr) \cup \bigl(\F_{\Reg(\epsilon)}(\G \oplus \epsilon) \setminus \Acc\bigr)$. If $\G$ is governed the first set is empty, so $\G \oplus \epsilon$ is governed exactly when the second set is empty, which is the admission condition. The second claim follows by induction on the number of events. For the third, if $\epsilon$ is admitted the second set is empty and the first is contained in $\F(\G) \setminus \Acc$.
\end{proof}
The check is thus \emph{exact}: an admitted event never breaks the invariant, and a held event always would. It does not depend on how the findings in the region are computed, which lets Section~\ref{sec:protocol} compute them on secret shares.

Owners act concurrently, so events must be admitted concurrently as well. Define the \emph{footprint} $\Phi(\epsilon)$ of an event as the vertices whose index entries its region reads or its application writes: $\key(e) \cup \bigcup_{r \in \key(e)} \mathrm{rd}(r)$ for data and usage rules, $\key(e) \cup \bigcup_{t \in \key(e)} \reads(d_1(t))$ for tool rules, and $\{t\} \cup \reads(d_1(t)) \cup \reads(d)$ for $\mathsf{redesc}(t, d)$ and $\mathsf{approve}(t)$.

\begin{proposition}[Commutation]\label{prop:commute}
If $\Phi(\epsilon_1) \cap \Phi(\epsilon_2) = \emptyset$, then the regions of $\epsilon_1$ and $\epsilon_2$ are disjoint and are the same in $\G$, $\G \oplus \epsilon_1$ and $\G \oplus \epsilon_2$, and both admission verdicts are the same in either order of application.
\end{proposition}
\begin{proof}
Every slot of $\Reg(\epsilon_i)$ contains a vertex of $\Phi(\epsilon_i)$ (a key vertex for pairs, a resource and a tool for cells), so disjoint footprints give disjoint regions. The region of $\epsilon_1$ is computed from the index entries of $\Phi(\epsilon_1)$, which $\epsilon_2$ does not write, and vice versa. By \Cref{thm:locality} applied to $\epsilon_2$, the findings of $\Reg(\epsilon_1)$ are the same in $\G \oplus \epsilon_1$ and $\G \oplus \epsilon_1 \oplus \epsilon_2$, so the verdict on $\epsilon_1$ does not depend on whether $\epsilon_2$ was applied first; symmetrically for $\epsilon_2$.
\end{proof}
Events with overlapping footprints are serialised, for instance by locking footprint vertices; in the AWS history of Section~\ref{sec:evaluation}, half of the changes touch at most \awsChangedMed{} action vertices, so conflicts between concurrent changes are rare.

\subsection{Algorithms}
Algorithm~\ref{alg:region} computes the region from two indexes: for each layer and key vertex $v$ the rules $\mathrm{Idx}_\ell[v]$ that constrain it, and for each resource the tools $\mathrm{rd}(r)$ that read it. Algorithm~\ref{alg:admit} evaluates the region in the successor state and decides. The findings of the region are kept in a store keyed by slot, so that \eqref{eq:locality} is applied by overwriting the region's entries.

\begin{algorithm}[t]
\caption{$\textsc{Region}(\epsilon)$: slots whose finding may change}\label{alg:region}
\begin{algorithmic}[1]
\Require event $\epsilon$; indexes $\mathrm{Idx}_\ell[v]$ and $\mathrm{rd}(r)$; population $S$
\Ensure candidate pairs $P$, cross-layer keys $X$, tools $T_u$ whose usage and drift cells must be re-evaluated
\State $P, X, T_u \gets \emptyset$
\If{$\epsilon = \mathsf{ins}(e)$ or $\mathsf{del}(e)$, with $e$ in layer $\ell$}
  \State $P \gets \{\{e, e'\} : v \in \key(e),\ e' \in \mathrm{Idx}_\ell[v] \setminus \{e\}\}$ \Comment{\Cref{lem:overlap,lem:subsume}}
  \If{$\ell = \mathsf{data}$} $X \gets \{(t, r) : r \in \key(e),\ t \in \mathrm{rd}(r)\}$
  \ElsIf{$\ell = \mathsf{tool}$} $X \gets \{(t, r) : t \in \key(e),\ r \in \reads(d_1(t))\}$
  \Else\ $T_u \gets \bigcup_{r \in \key(e)} \mathrm{rd}(r)$
  \EndIf
\ElsIf{$\epsilon = \mathsf{redesc}(t, d)$}
  \State $X \gets \{(t, r) : r \in \reads(d_1(t)) \cup \reads(d)\}$;\quad $T_u \gets \{t\}$
\ElsIf{$\epsilon = \mathsf{approve}(t)$}
  \State $T_u \gets \{t\}$
\EndIf
\State \Return $P$, $X \times S$ (cells $(u, t, r)$), $T_u$
\end{algorithmic}
\end{algorithm}

\begin{algorithm}[t]
\caption{$\textsc{Admit}(\epsilon)$: exact admission of one event}\label{alg:admit}
\begin{algorithmic}[1]
\Require governed state $\G$ with finding store $\mathcal{S}$ (slot $\mapsto$ finding); accepted set $\Acc$
\Ensure verdict, and findings for their owners
\State $(P, X, T_u) \gets \textsc{Region}(\epsilon)$ \Comment{Algorithm~\ref{alg:region}}
\State $\G' \gets \G \oplus \epsilon$ \Comment{tentative; indexes updated on a copy of the footprint}
\State $N \gets \emptyset$
\ForAll{$\{e, e'\} \in P$}
  \State $N \gets N \cup \textsc{PairCheck}(e, e')$ \Comment{local if one owner, else Algorithm~\ref{alg:pairsecure}}
\EndFor
\ForAll{$(t, r) \in X$}
  \State $N \gets N \cup \textsc{CellCheck}(t, r, S)$ \Comment{local if one owner, else Algorithm~\ref{alg:ldsc}}
\EndFor
\ForAll{$t \in T_u$}
  \State $N \gets N \cup \Viol(d'_1(t)) \cup \bigl(\Viol(d'_1(t)) \setminus \Viol(d'_0(t))\bigr)$ \Comment{evaluated by the owners of $\reads(t)$}
\EndFor
\If{$N \subseteq \Acc$}
  \State commit $\G'$; overwrite the entries of $P$, $X \times S$ and $T_u$ in $\mathcal{S}$ with $N$ \Comment{\eqref{eq:locality}}
  \State \Return \textsc{accept}
\Else
  \State \Return \textsc{hold}, and send each finding of $N \setminus \Acc$ to the owners it involves, ranked by Section~\ref{sec:ranking}
\EndIf
\end{algorithmic}
\end{algorithm}

\begin{proposition}[Cost]\label{prop:cost}
Let $n_v = |\mathrm{Idx}_\ell[v]|$, $m = |A_\ell|$ and $N_\ell$ the number of rules of layer $\ell$. Admitting a rule event costs $O\bigl(m \sum_{v \in \key(e)} n_v\bigr)$ pair tests and $O\bigl(|S|\,(m \sum_{v \in \key(e)} n_v + |X|)\bigr)$ for the cells, where $|X| \le |\key(e)| \cdot \max_r |\mathrm{rd}(r)|$ for data rules and $|\key(e)| \cdot \max_t |\reads(t)|$ for tool rules. Full recomputation costs $O(m \sum_v n_v^2)$ pair tests and $O\bigl(|S|(m N_\ell + \sum_t |\reads(t)|)\bigr)$ for the cells.
\end{proposition}
\begin{proof}
Each pair test is a per-attribute set comparison (\Cref{lem:overlap,lem:subsume}). A decision row $\dec_\ell(\cdot, v)$ over $S$ is recomputed from the $n_v$ rules on $v$ in $O(|S| m n_v)$, and each cell key then costs $O(|S|)$. Full recomputation does the same for every vertex and every pair inside each bucket.
\end{proof}
When rules spread over many key vertices, $n_v$ is small and the ratio between full recomputation and admission grows with the number of vertices; Section~\ref{sec:scale} measures it.

\begin{example}[Running example, admission]\label{ex:admission}
Before the three events, both clinicians and analysts form blocked grants on \texttt{summarise} and \textit{ehr}, since $C$ denies $V$'s service account; suppose $C$ and $V$ had accepted them.
\textbf{$\epsilon_1 = \mathsf{ins}(e_5)$} has $\key(e_5) = \{\textit{ehr}\}$ and $\mathrm{rd}(\textit{ehr}) = \{\texttt{summarise}\}$, so its region is the pair $\{e_5, e_3\}$ (both of $C$, checked locally) and the cells $(u, \texttt{summarise}, \textit{ehr})$. In the successor state the two blocked grants disappear and $(\textit{analyst}, \texttt{summarise}, \textit{ehr})$ becomes an amplification, as Theorem~\ref{thm:duality} predicts. The event is held, and the finding goes to $C$ and $V$.
\textbf{$\epsilon_2 = \mathsf{redesc}(\texttt{export}, d)$} has the cells of \texttt{export} as its region. The usage cell $(\texttt{export}, \textit{claims}, R/\mathrm{US})$ is denied by $C$'s usage rule and absent under the approved description: a usage violation and a drift, sent to $C$ and $V$.
\textbf{$\epsilon_3 = \mathsf{ins}(e_7)$}, a tool rule on \texttt{export}, has the cells $(u, \texttt{export}, \textit{claims})$ as its region; interns are offered \texttt{export} but $C$ denies them \textit{claims}: a blocked grant. $C$ may accept it, since the call will simply fail, or $V$ may withdraw the rule.
Figure~\ref{fig:hg}b shows the three new hyperedges and their key vertices. None of the three checks reads any rule of the hypergraph outside the region of its event.
\end{example}
